\section{What localized Fourier--Laplace forgets}
\label{sec:fourier-boundary}

The incomplete-Gamma jump has an elementary differential-algebraic shadow.
For \(g(x)=\Gamma(1-\alpha,x)\),
\[
 (\partial_x+1+\alpha/x)\partial_xg=0.                         \tag{6.1}
\]
Differentiation forgets the constant solution, and that constant is the
additive datum in the continuation formula.

Let \(D_t=\C\langle t,\partial_t\rangle\).  Under algebraic Fourier transform,
\[
 \C[t]=D_t/D_t\partial_t
 \quad\longmapsto\quad
 \delta_0=D_\tau/D_\tau\tau.                                \tag{6.2}
\]
Localizing at \(\tau\) kills \(\delta_0\).  Localized partial
Fourier--Laplace therefore cannot recover the extension constant in (6.1).
The four-term GKZ/Gauss--Manin comparison in \cite{RSSW} uses the same
nonconservativity.

\begin{proposition}[Punctual Fourier corner]
\label{prop:punctual-corner}
Let
\[
 \delta_{00}=
 D_{\tau_1,\tau_2}/D_{\tau_1,\tau_2}(\tau_1,\tau_2).
\]
Localization at either \(\tau_i\) kills \(\delta_{00}\), whereas
\[
 \FL^{-1}(\delta_{00})=\OO_{\A^2}                            \tag{6.3}
\]
up to shift and sign conventions.  Thus an object supported at the
intersection of two deleted Fourier boundaries can have generic-rank-one
inverse transform.
\end{proposition}

\begin{proof}
Use the Weyl-algebra Fourier automorphism
\(x_i\mapsto-\partial_{\tau_i}\) and
\(\partial_{x_i}\mapsto\tau_i\).  It exchanges
\(D_x/D_x(\partial_{x_1},\partial_{x_2})\) with \(\delta_{00}\).
Since \(\tau_i\) acts nilpotently on the latter, inverting either coordinate
kills it.  The inverse transform is the constant module, which has generic
rank one.
\end{proof}

Tensoring Proposition~\ref{prop:punctual-corner} with the
\(\zetaSix\) rank-one connection, its hyperbolic dual, and
\(\Q[N]/(N^3)\) preserves every conclusion.  Bare double-boundary support is
therefore not a rank-zero condition.

Nor does marking one punctual vector make it unique.  Let \(C\) be any
finite-dimensional coefficient space and put
\(\mathcal B=\delta_{00}\otimes C\).  For
\(G\in\operatorname{GL}(C)\), two receivers can differ by \(G\) while
having the same support and characteristic cycle.  After hyperbolic doubling,
\(G\oplus G^{-t}\) preserves a perfect pairing.  It can still change a chosen
point row.  Consequently, a claim that the punctual quotient is generated by
one marked line is additional geometry.  Enlarging the discarded subspace to
the kernel of the point row makes the quotient one-dimensional only
tautologically: invariance of that kernel is the theorem one seeks.

The unlocalized boundary map
\[
 \FL(M)\longrightarrow j_*j^*\FL(M),\qquad
 j:\Gm\hookrightarrow\A^1,                                  \tag{6.4}
\]
retains the missing object.  In a two-coordinate problem one must also retain
the orientation, duality, and \(!\to *\) or
\(\operatorname{can}/\operatorname{var}\) maps of the complete Cech boundary
square.  Proposition~\ref{prop:punctual-corner} shows that retaining the
object is necessary, not that its support alone proves the desired vanishing.
